%==============================================================================
% Two-page IEEE proposal in the "Jahin" defensive-writing style.
% Compile: pdflatex proposal_jahin && pdflatex proposal_jahin
%==============================================================================
\documentclass[conference]{IEEEtran}
\IEEEoverridecommandlockouts
\usepackage{cite}
\usepackage{amsmath,amssymb}
\usepackage{booktabs}
\usepackage[T1]{fontenc}
\usepackage{enumitem}
\setlist[itemize]{leftmargin=1.1em,noitemsep,topsep=1pt}
\setlength{\textfloatsep}{5pt}
\newtheorem{definition}{Definition}

\begin{document}

\title{A Code-Based Approach to Secure Aggregation\\ in Federated Learning}

\author{
\IEEEauthorblockN{Md Wahiduzzaman Suva (26-94088-2)\qquad Esm-e Moula Chowdhury Abha (26-94089-2)}
\IEEEauthorblockA{Department of Computer Science, American International University-Bangladesh\\
wshuvo360@gmail.com}
\thanks{Two-page research proposal. Design and evaluation are planned work; the reported
implementation is a preliminary prototype.}
}

\maketitle

\begin{abstract}
Federated learning (FL) trains shared models without centralizing data, yet the model updates
clients exchange leak private information, and the secure-aggregation protocols that prevent this
rely on key establishment that a quantum computer would break, exposing recorded updates under
``harvest-now, decrypt-later.'' Post-quantum FL secures this layer almost exclusively with
\emph{lattice} cryptography, concentrating an entire ecosystem on one hardness assumption that a
single cryptanalytic advance would defeat. We propose CB-SAFE, a KEM-agnostic secure-aggregation
framework whose post-quantum confidentiality is established by the NIST-selected code-based KEM
HQC rather than a lattice KEM, providing the same NIST security levels under an independent
(syndrome-decoding) assumption. CB-SAFE exposes the KEM as a pluggable module (Definition~1), runs
code-based HQC or lattice Kyber/ML-KEM under one masking protocol, and applies cross-round key
amortization, ciphertext batching, and compression so the larger code-based payloads stay
practical. A prototype liboqs implementation establishes initial feasibility, and we outline an
equal-security evaluation against the lattice baseline across three public datasets (FEMNIST,
CIFAR-10, Edge-IIoTset) and five metric families. To our knowledge, this is the first FL
secure-aggregation framework built on a code-based KEM; the implementation will be released
open-source.
\end{abstract}

\begin{IEEEkeywords}
Post-quantum cryptography, federated learning, secure aggregation, HQC, code-based cryptography,
cryptographic diversity.
\end{IEEEkeywords}

\section{Introduction}
Federated learning (FL) lets many clients train a shared model under a coordinating server while
keeping their data local, but the transmitted model updates are not private: gradient-inversion
and membership-inference attacks reconstruct training samples from them. Secure aggregation removes
this leakage by ensuring the server observes only the sum of updates, through pairwise additive
masks that cancel on aggregation. The confidentiality of these masks rests on classical key
exchange, which Shor's algorithm renders insecure once a cryptographically relevant quantum
computer exists, exposing recorded updates retroactively under ``harvest-now, decrypt-later.''

The natural defense replaces the key exchange with a post-quantum KEM, and recent post-quantum FL
does this almost always with lattice cryptography: lattice multi-stage secret sharing
\cite{pqsf2024}, lattice key-homomorphic pseudo-random functions \cite{qinxu2025}, and lattice
signature aggregation \cite{narsimhulu2026}, with non-lattice work rare and code-free
\cite{singh2025pqfog} and surveys confirming lattice dominance \cite{nguyen2025lattice}.
Concentrating confidentiality on one hardness assumption across an ecosystem is a systemic risk:
a single cryptanalytic advance would break every such system at once.

NIST anticipated this risk when, in March~2025, it selected the code-based KEM HQC (whose security
reduces to syndrome decoding of quasi-cyclic codes rather than to lattice problems) as a
diversified backup to ML-KEM \cite{nistir8545}. The same diversification belongs inside FL's
confidentiality layer; we realize it by making the KEM a pluggable module, so code-based and
lattice key establishment run under one protocol and differ only in the swapped primitive. Our
contributions are:

\begin{itemize}
\item \textbf{A KEM-agnostic secure-aggregation protocol} that formalizes masked aggregation over
any IND-CCA2 KEM (Definition~1), making the post-quantum primitive a single replaceable module.
\item \textbf{The first code-based instantiation, scope-qualified:} to our knowledge, CB-SAFE is
the first FL secure-aggregation framework whose post-quantum confidentiality is established by the
NIST-selected code-based KEM (HQC) rather than a lattice KEM, at the same NIST security levels.
\item \textbf{Ciphertext-handling for practicality:} cross-round key amortization, ciphertext
batching, and compression that keep the code-based payloads practical inside the FL loop.
\item \textbf{A planned reproducible implementation} (liboqs, Flower, PyTorch), prototyped to
establish feasibility, with an equal-security evaluation across three datasets and five metric
families (Section~IV).
\end{itemize}

\section{Related Work}
Post-quantum FL secures aggregation almost exclusively with lattices: multi-stage secret sharing
\cite{pqsf2024}, key-homomorphic PRFs \cite{qinxu2025}, and signature aggregation
\cite{narsimhulu2026}, each advancing efficiency within the lattice family but inheriting its
single-assumption exposure; non-lattice work is scarce and code-free (e.g., multivariate
cryptography with blockchain \cite{singh2025pqfog}), and surveys confirm lattice dominance
\cite{nguyen2025lattice}.

Three commitments are common to the work above and motivate CB-SAFE's departure.
(i)~Confidentiality rests on a single lattice assumption \cite{pqsf2024,qinxu2025,narsimhulu2026};
CB-SAFE establishes it from an independent code-based assumption selected by NIST for exactly this
diversity \cite{nistir8545}. (ii)~The post-quantum primitive is hard-wired into each protocol;
CB-SAFE exposes it behind a KEM-agnostic interface (Definition~1) so the code-based and lattice
variants differ only in one module. (iii)~No prior framework brings a standard-track code-based KEM
into FL secure aggregation. To our knowledge, none satisfies any of these three.

\section{Proposed Framework: CB-SAFE}
\textbf{System and threat model.} We consider cross-device/cross-silo FL with a server and $N$
clients running FedAvg over non-IID partitions. The server is honest-but-curious and may collude
with up to $t<N/2$ clients; we additionally model malicious clients submitting poisoned updates and
clients dropping out mid-round. The goal is that no coalition below the threshold learns any
individual update $x_i$, only the aggregate $\sum_i x_i$.

\begin{definition}[KEM-agnostic masked aggregation]
Let $\mathcal{K}=(\mathrm{KeyGen},\mathrm{Encaps},\mathrm{Decaps})$ be an IND-CCA2 KEM and
$H:\{0,1\}^*\!\to\mathcal{R}$ a key-derivation function into the update space $\mathcal{R}$. For an
ordered pair $(i,j)$, client $i$ computes $(c_{ij},ss_{ij})=\mathrm{Encaps}(pk_j)$ and the mask
$m_{ij}=H(ss_{ij})$, and client $j$ recovers $ss_{ij}=\mathrm{Decaps}(sk_j,c_{ij})$. Client $i$
submits $y_i = x_i + \sum_{j>i} m_{ij} - \sum_{j<i} m_{ij}$; the server returns
$\sum_i y_i = \sum_i x_i$, since each $m_{ij}$ appears once with each sign.
\end{definition}
Instantiating $\mathcal{K}$ with HQC yields the code-based variant; with Kyber/ML-KEM
\cite{fips203}, the lattice baseline (Table~I). As the protocol uses $\mathcal{K}$ only through its
IND-CCA2 interface, correctness and privacy are identical across instantiations. Dropout is handled
by secret-sharing each client's seed so the server can cancel the masks of non-responding clients.

\textbf{Making HQC practical.} Replacing the KEM is not a drop-in substitution: code-based KEMs use
larger ciphertexts and key material than lattice KEMs, so naive per-round encapsulation would
inflate communication. CB-SAFE keeps the code-based variant practical by amortizing key
establishment across rounds (long-lived HQC public keys are reused and only masks refreshed, so the
$O(N)$ ciphertext exchange becomes a setup cost) and by batching each client's pairwise ciphertexts
and compressing key material. HQC cannot sign and no code-based signature is standardized, so we
confine the code-based primitive to confidentiality and, where authentication is needed, compose a
standardized lattice signature (ML-DSA), with only the KEM layer code-based.

\begin{table}[t]
\caption{Configurations compared under the KEM-agnostic CB-SAFE protocol (Definition~1), at
identical NIST security levels; CB-SAFE adds the code-based variant absent from prior
post-quantum FL.}
\label{tab:config}
\centering
\footnotesize
\setlength{\tabcolsep}{4pt}
\begin{tabular}{@{}llll c@{}}
\toprule
Configuration & Family & Assumption & PQ-safe & Level\\
\midrule
Plain FedAvg          & none       & n/a               & No  & n/a\\
CB-SAFE / lattice     & lattice    & Module-LWE        & Yes & 1/3/5\\
CB-SAFE / code-based  & code-based & Syndrome decoding & Yes & 1/3/5\\
\bottomrule
\end{tabular}
\end{table}

\section{Evaluation Plan and Metrics}
We compare three configurations, plain FedAvg, the lattice (Kyber/ML-KEM) baseline, and the
code-based (HQC) CB-SAFE variant (Table~I), across three public datasets in two domains: FEMNIST
and non-IID CIFAR-10 (image classification) and Edge-IIoTset (IoT intrusion detection)
\cite{narsimhulu2026,singh2025pqfog}, all at matched NIST levels (1/3/5). We report five metric
families: (1)~\emph{communication}, bytes/round/client separating model payload from key and
ciphertext overhead; (2)~\emph{computation}, keygen/encaps/decaps time and per-round wall-clock;
(3)~\emph{utility}, accuracy vs.\ rounds, which must coincide with the plain baseline since correct
masking preserves the aggregate (Definition~1); (4)~\emph{robustness}, accuracy under
malicious/colluding-client fractions and dropout; and (5)~\emph{security level}. Utility and
robustness use paired tests with multiple-comparison correction over repeated seeds, establishing
the first code-based reference point for post-quantum secure aggregation in FL.

\section{Conclusion}
Post-quantum FL has equated ``quantum-safe'' with ``lattice-based,'' leaving privacy-critical
deployments exposed to a single cryptanalytic failure. CB-SAFE reframes the post-quantum KEM as a
replaceable module and supplies the first code-based occupant of that slot, bringing the
cryptographic diversity for which NIST standardized HQC into federated learning at equal security
levels. By exposing the KEM behind a KEM-agnostic interface, CB-SAFE makes cryptographic agility a
configuration choice as the post-quantum landscape evolves.

\begin{thebibliography}{1}
\normalsize
\bibitem{nistir8545} National Institute of Standards and Technology, ``Status report on the fourth
round of the NIST post-quantum cryptography standardization process,'' NIST IR 8545, Mar.\ 2025.
\bibitem{fips203} National Institute of Standards and Technology, ``Module-lattice-based
key-encapsulation mechanism standard,'' FIPS 203, Aug.\ 2024.
\bibitem{pqsf2024} X.\ Zhang, H.\ Deng, R.\ Wu, J.\ Ren, and Y.\ Ren, ``PQSF: Post-quantum secure
privacy-preserving federated learning,'' \emph{Sci.\ Rep.}, vol.\ 14, Art.\ no.\ 23553, 2024.
\bibitem{qinxu2025} X.\ Qin and R.\ Xu, ``Efficient post-quantum cross-silo federated learning based
on key homomorphic pseudo-random function,'' \emph{Mathematics}, vol.\ 13, no.\ 9, Art.\ no.\ 1404,
2025.
\bibitem{narsimhulu2026} P.\ Narsimhulu, P.\ Chithaluru, and R.\ Aluvalu, ``Efficient post-quantum
cryptographic signature aggregation for low-latency distributed networks,'' \emph{EURASIP J.\ Inf.\
Secur.}, vol.\ 2026, Art.\ no.\ 6, 2026.
\bibitem{singh2025pqfog} J.\ Singh, P.\ Singh, A.\ Kaur, and M.\ Hedabou, ``Post-quantum secure
fog-edge computing using federated learning with blockchain,'' \emph{J.\ Supercomput.}, vol.\ 81,
Art.\ no.\ 1249, 2025.
\bibitem{nguyen2025lattice} H.\ Nguyen, S.\ Huda, Y.\ Nogami, and T.\ T.\ Nguyen, ``Security in
post-quantum era: A comprehensive survey on lattice-based algorithms,'' \emph{IEEE Access}, vol.\
13, pp.\ 89003--89024, 2025.
\end{thebibliography}

\end{document}
